\chapter{Obstacle existence, comparison, and reaction}\label{sec:obstacle-variational}
We use the admissible set and the sign convention of Definition \ref{def:obstacle-problem}: the energy is $\mathcal E_f(v)=\frac12\int_U|Dv|^2-\int_Ufv$, with a lower obstacle. The direct method and one-sided variations already yield existence, comparison, and a reaction measure. The zero-forcing version is developed in \cite{Fig18}, Sections 2--3, pp.~2--5; the proofs below retain an $L^2$ forcing term.

\begin{theorem}[Existence and uniqueness for the scalar obstacle problem]\label{thm:obstacle-existence}
Under the hypotheses of Definition \ref{def:obstacle-problem}, there is a unique minimizer $u\in K$ of $\mathcal E_f$. Uniqueness is as an $H^1$ function, hence up to almost-everywhere equality. No regularity of its free boundary is part of this assertion.
\end{theorem}
\begin{proof}
Fix $v_0\in K$ and write $v=v_0+w$, $w\in H^1_0(U)$. Poincare and Cauchy--Schwarz give
\[
 \mathcal E_f(v)\geq\tfrac14\|Dw\|_2^2-C,
\]
where $C$ depends on $v_0,f,U$. Thus a minimizing sequence is bounded in $H^1$. Pass to a weakly convergent subsequence. The affine space $g+H^1_0(U)$ is weakly closed. The obstacle constraint is also weakly closed: for every nonnegative $\eta\in L^2(U)$, the inequalities $\int(v_j-\psi)\eta\geq0$ pass to the limit and imply $v\geq\psi$ almost everywhere. The squared gradient norm is weakly lower semicontinuous and the forcing term is weakly continuous, so the limit minimizes.

If $u_1,u_2$ minimize, their midpoint is admissible. The parallelogram identity shows that its energy is the average energy minus $\frac18\|D(u_1-u_2)\|_2^2$. Thus the gradients agree. Since $u_1-u_2\in H^1_0(U)$, Poincare gives equality of the functions. This extends the proof of \cite{Fig18}, Proposition 2.1, to the stated forcing class.
\end{proof}

\begin{proposition}[The energy gap controls distance to the minimizer]\label{prop:obstacle-energy-gap}
For the minimizer $u$ and any $v\in K$,
\[
 \mathcal E_f(v)-\mathcal E_f(u)\geq\frac12\|D(v-u)\|_{L^2(U)}^2.
\]
In particular every minimizing sequence converges strongly to $u$ in $H^1(U)$.
\end{proposition}
\begin{proof}
Expand the quadratic energy as in Proposition \ref{prop:obstacle-inequality}. Its linear part is nonnegative by the variational inequality, leaving the displayed bound. Poincare controls the $L^2$ norm of $v-u$, since the boundary traces agree. Thus convergence of energies to the minimum implies strong convergence without taking a further subsequence.
\end{proof}

\begin{proposition}[Stability with respect to forcing]\label{prop:obstacle-forcing-stability}
Keep $U,g,\psi$ fixed and let $u_i$ minimize with forcing $f_i\in L^2(U)$. Then
\[
 \|D(u_1-u_2)\|_2\leq\|f_1-f_2\|_{H^{-1}(U)},
\]
where the dual norm uses $\|D\phi\|_2$ on $H^1_0(U)$.
\end{proposition}
\begin{proof}
Use $u_2$ as the competitor for $u_1$ and $u_1$ for $u_2$ in the two variational inequalities. After addition and rearrangement,
\[
 \|D(u_1-u_2)\|_2^2\leq\int_U(f_1-f_2)(u_1-u_2).
\]
The difference belongs to $H^1_0(U)$. Apply the dual-norm bound and divide by its gradient norm when this norm is nonzero. The zero case is immediate.
\end{proof}

\begin{proposition}[Comparison for ordered data]\label{prop:obstacle-comparison}
On the same bounded Lipschitz domain, suppose $\psi_1\leq\psi_2$ almost everywhere, the boundary traces satisfy $\operatorname{Tr}g_1\leq\operatorname{Tr}g_2$, and $f_1\leq f_2$ almost everywhere. Assume both admissible classes are nonempty. Then their minimizers satisfy $u_1\leq u_2$ almost everywhere.
\end{proposition}
\begin{proof}
Sobolev truncation and the ordered traces give $w=(u_1-u_2)^+\in H^1_0(U)$. The functions $u_1-w=\min\{u_1,u_2\}$ and $u_2+w=\max\{u_1,u_2\}$ are admissible in the first and second problems respectively. The obstacle ordering is needed for the first assertion. Their variational inequalities imply
\[
 \int_U|Dw|^2
 =\int_U(Du_1-Du_2)\cdot Dw
 \leq\int_U(f_1-f_2)w\leq0.
\]
Poincare gives $w=0$. The Sobolev positive-part chain rule and compatibility of traces with Lipschitz truncation are the functional-analytic inputs, with background in \cite{SimonPDE}, Chapters 5--9.
\end{proof}

\begin{proposition}[The nonnegative reaction measure]\label{prop:obstacle-reaction}
For a minimizer $u$, the distribution
\[
 \lambda=-\Delta u-f
\]
is a nonnegative Radon measure in $U$. In particular $-\Delta u\geq f$ distributionally. The assertion is local finiteness on compact subsets, not necessarily finite mass on all of $U$.
\end{proposition}
\begin{proof}
If $\phi\in C_c^\infty(U)$ is nonnegative, $u+t\phi$ is admissible for $t\geq0$. The variational inequality gives $\langle\lambda,\phi\rangle=\int Du\cdot D\phi-\int f\phi\geq0$. A positive distribution has order zero locally: for test functions supported in a compact $K$, choose a smooth nonnegative cutoff $\eta$ equal to one near $K$. Positivity applied to $\|\phi\|_\infty\eta\pm\phi$ bounds $|\langle\lambda,\phi\rangle|$ by $\|\phi\|_\infty\langle\lambda,\eta\rangle$. Extension to continuous test functions and the Riesz representation theorem give a nonnegative Radon measure. This is the forced version of \cite{Fig18}, Proposition 3.1.
\end{proof}

\begin{proposition}[The equation away from contact]\label{prop:obstacle-noncontact-equation}
Suppose additionally that $u$ and $\psi$ have continuous representatives. On the open set $O=\{u>\psi\}$,
\[
 -\Delta u=f,
\]
and the reaction measure vanishes. Consequently its support in $U$ is contained in the contact set $\{u=\psi\}$.
\end{proposition}
\begin{proof}
A smooth test function compactly supported in $O$ has support on which $u-\psi$ has a positive minimum. Thus both $u+t\phi$ and $u-t\phi$ are admissible for sufficiently small $t>0$. Their variational inequalities give equality in the weak equation. Since the reaction vanishes on every such test function, it vanishes as a measure on $O$. The contact set is relatively closed by continuity. Compare \cite{Fig18}, Corollary 3.4. Without a continuous representative, a statement involving pointwise contact sets requires additional care.
\end{proof}

\begin{proposition}[The minimizer is the least admissible supersolution]\label{prop:obstacle-least-supersolution}
If $v\in K$ satisfies
\[
 \int_U Dv\cdot D\phi\geq\int_U f\phi
 \quad\text{for every nonnegative }\phi\in H^1_0(U),
\]
then the minimizer $u$ satisfies $u\leq v$ almost everywhere.
\end{proposition}
\begin{proof}
The variational inequality with $u+\phi$ shows that $u$ itself is an admissible supersolution. Set $w=(u-v)^+\in H^1_0(U)$. The competitor $\min\{u,v\}=u-w$ lies in $K$, so the variational inequality gives $\int Du\cdot Dw\leq\int fw$. The supersolution inequality gives the reverse bound for $Dv$. Subtraction yields $\int|Dw|^2\leq0$, and Poincare gives the result. This is an order characterization within the stated common boundary class, not an assertion about arbitrary boundary data.
\end{proof}

\begin{proposition}[Subtracting the obstacle and fixing the forcing sign]\label{prop:obstacle-normalization}
Assume $\psi\in C^2(U)$ and put $v=u-\psi$. Locally the new obstacle is zero, and differences of energies between competitors with the same boundary values are those of
\[
 \mathcal J_a(v)=\int\left(\tfrac12|Dv|^2+av\right),
 \qquad a=-f-\Delta\psi.
\]
On the noncontact region, whenever Proposition \ref{prop:obstacle-noncontact-equation} applies, $\Delta v=a$. The normalized model $a=1$ therefore has a positive linear term in its energy.
\end{proposition}
\begin{proof}
Expand $|D(\psi+v)|^2$ and compare two competitors whose difference $w$ lies in $H^1_0$ on a relatively compact subdomain. The cross term is $\int D\psi\cdot Dw=-\int(\Delta\psi)w$. Together with $-\int fw$ this gives $\int aw$. The remaining terms depending only on $\psi$ cancel. The equation follows by subtracting $\Delta\psi$ from $\Delta u=-f$. This is the sign convention behind \cite{Fig18}, Section 5; that source begins with $f=0$ and chooses $-\Delta\psi=1$.
\end{proof}
